\documentclass[10pt,a4paper]{amsart}
\usepackage[margin=19mm]{geometry}
\usepackage{amsmath,amssymb,mathtools}
\usepackage[scr=rsfs,cal=euler]{mathalfa}
\usepackage{xfrac}
\usepackage[unicode,hidelinks,bookmarks=false]{hyperref}
\newcommand{\ie}{i.e.\ }
\newcommand{\cf}{cf.\ }
\newcommand{\resp}{respectively\ }
\newcommand{\Subsection}{\subsection}
\providecommand{\nobreakdash}{\nobreak\hbox{-}}
\setlength{\emergencystretch}{2em}
\multlinegap0pt
\usepackage{amssymb}
\newtheorem{thm}{Theorem}
\title{Nuclear Weighted composition operators between different $L^p$-spaces}
\author{S. Al Ghafri, Y. Estaremi, S. Shamsigamchi}
\date{Source: arXiv:2603.20681v1 (CC BY 4.0)}
\begin{document}
\maketitle

\noindent\emph{Source and licence.} Nuclear Weighted composition operators between different $L^p$-spaces. Version arXiv:2603.20681v1 (CC BY 4.0), distributed by arXiv under the Creative Commons Attribution 4.0 International licence, http://creativecommons.org/licenses/by/4.0/, as recorded in the arXiv licence metadata for this exact version and shown on the abstract page https://arxiv.org/abs/2603.20681v1. Full licence notice: \texttt{licenses/arxiv-2603.20681-CC-BY-4.0.txt}.\par\smallskip

\noindent\textit{Editorial scope.} Counted unit: the article's thm n.pp (source0.tex lines 182--184). The article's complete original proof of it is delivered with it (source0.tex lines 185--219, one \texttt{proof} environment of 35 lines). No mathematics is written, repaired or re-derived: the statement and the proof are the article's own lines, in the article's own notation, and no retained line is substituted or re-flowed.  Retained uncounted as the source-local context the proof needs: lines 167--179.  Nothing was omitted from any retained span, so the package carries no omission record.  The retained lines cite 1 works, whose entries are transcribed verbatim from the article's own bibliography source in the e-print and delivered with short editorial labels recorded in the item's reference mapping.  No retained line contains a figure environment, an image-inclusion command or drawing code, so no image is delivered and the package declares no asset dependency.  Editorial formatting only, in addition: an AMS \texttt{amsart} preamble that restates the article's own declarations and macros used by the retained lines, namely the environments \texttt{thm} on separate counters, together with the standard packages those lines need.  The retained environments print in this document as Theorem 1 in order of appearance, whereas the article prints the same items with the numbering of its own full document.  The complete native source of the article as distributed in the arXiv e-print for this exact version is bundled unchanged as \texttt{sources/196870/source0.tex}.
\begin{thm}[Pietsch domination theorem {\cite[Theorem~2.12]{dj}}]\label{thm:Pietsch}
Let $1\le p<\infty$ and let $T:X\to Y$ be absolutely $p$--summing.
Then there exist $C>0$ and a Borel probability measure $\mu$ on the weak$^{*}$--compact unit ball $B_{X^*}$ such that
\[
\|Tx\|
\le C\left(\int_{B_{X^*}} |x^*(x)|^p\,d\mu(x^*)\right)^{1/p},
\qquad x\in X.
\]
In particular, for $p=1$,
\[
\|Tx\|\le C\int_{B_{X^*}} |x^*(x)|\,d\mu(x^*).
\]
\end{thm}

\begin{thm}\label{n.pp}
Let $(X, \mu, \Sigma)$ is a $\sigma$-finite measure spaces and $X=\cup^{\infty}_{n=1}A_n\cup B$, in which $A_n$'s are disjoint $\sigma$-atoms and $B$ is the non-atomic part. If $1\leq p<\infty$  and $W=uC_{\varphi}:L^p(\mu)\rightarrow L^p(\mu)$ is the weighted composition operator, then $W$ is nuclear if and only if $\mu(\varphi^{-1}(B)\cap S(u))=0$ and $\sum^{\infty}_{n=1}J_p(A_n)^{\frac{1}{p}}<\infty$.
\end{thm}

\begin{proof}
By assumptions we have $X=\cup^{\infty}_{n=1}A_n\cup B=\cup^{\infty}_{n=1}\varphi^{-1}(A_n)\cup \varphi^{-1}(B)$. Let $\mu(\varphi^{-1}(B)\cap S(u))=0$ and $\sum^{\infty}_{n=1}J_p(A_n)^{\frac{1}{p}}<\infty$. Then for every $f\in L^p(\mu)$,
\begin{align*}
u.f\circ \varphi&=u.\sum^{\infty}_{n=1}f\circ\varphi.\chi_{\varphi^{-1}(A_n)}\\
&=u.\sum^{\infty}_{n=1}(f.\chi_{A_n})\circ\varphi\\
&=u.\sum^{\infty}_{n=1}f(A_n).\chi_{A_n}\circ\varphi\\
&=\sum^{\infty}_{n=1}\frac{u.\chi_{\varphi^{-1}(A_n)}}{\mu(A_n)}\int_{A_n}fd\mu\\
&=\sum^{\infty}_{n=1}\phi_n(f)g_n,
\end{align*}
in which $g_n=\frac{u.\chi_{\varphi^{-1}(A_n)}}{\mu(A_n)}$ and $\phi_n(f)=\int_{A_n}fd\mu$. It is easy to see that $g_n\in L^p(\mu)$, $\phi_n\in (L^p(\mu))^*$ and $\|g_n\|_p=\frac{(J_p(A_n))^{\frac{1}{p}}}{\mu(A_n)^{\frac{1}{p'}}}$, $\|\phi_n\|\leq\mu(A_n)^{\frac{1}{p'}}$ (for $1<p<\infty$) and also $\|g_n\|_1=J_1(A_n)$ and $\|\phi_n\|\leq1$ (for $p=1$), in which $\frac{1}{p}+\frac{1}{p'}=1$. This implies that 
$$\sum^{\infty}_{n=1}\|\phi_n\|.\|g_n\|_p\leq \sum^{\infty}_{n=1}J_p(A_n)^{\frac{1}{p}}<\infty.$$
So $W=u.C_{\varphi}$ is nuclear.\\
Conversely, let $W=u.C_{\varphi}$ be nuclear. Then it is absolutely summing and compact, so by Theorem \ref{thm:Pietsch}, there is a Borel probability measure $\nu$ on the weak$^*$ compact closed unit ball $B_{p'}$ of $(L^p(\mu))^*\simeq L^{p'}(\mu)$ (case $1<p<\infty$), $B_{\infty}$ of $(L^1(\mu))^*\simeq L^{\infty}(\mu)$ (case $p=1$) and $C>0$ such that
$$\|u.C_{\varphi}f\|_p\leq C \int_{B_{p'}}|\xi(f)|\nu(\xi), \ \ \ \ \ \ f\in L^p(\mu),$$
and 
$$\|u.C_{\varphi}f\|_1\leq C \int_{B_{\infty}}|\xi(f)|\nu(\xi), \ \ \ \ \ \ f\in L^1(\mu).$$
For each $n\in \mathbb{N}$, we set $f_n=\frac{\chi_{A_n}}{\mu(A_n)^{\frac{1}{p}}}$ and then we have
\begin{align*}
J_p(A_n)^{\frac{1}{p}}&=\frac{1}{\mu(A_n)^{\frac{1}{p}}}\left(\int_X|u|^p\chi_{A_n}\circ\varphi d\mu\right)^{\frac{1}{p}}\\
&=\|u.f_n\circ\varphi\|_p\\
&\leq C\int_{B_{p'}}|\xi(f_n)|\nu(\xi)\\
&=C\int_{B_{p'}}|\int_X f_n.g_{\xi}d\mu\nu(g_{\xi})|\\
&\leq C \int_{B_{p'}}\|f_n\|_p.\|g_{\xi}.\chi_{A_n}\|_{p'}\nu(g_{\xi})\\
&\leq C\int_{B_{p'}}\|g_{\xi}.\chi_{A_n}\|_{p'}\nu(g_{\xi}),\\
\end{align*}
and similarly $J_1(A_n)\leq C\int_{B_{\infty}}\|g_{\xi}.\chi_{A_n}\|_{\infty}\nu(g_{\xi})$. On the other hands compactness of $W$ implies that $\mu(\varphi^{-1}(B)\cap S(u))=0$. So $\|g_{\xi}\|_{p'}=\sum^{\infty}_{n=1}\|g_{\xi}.\chi_{A_n}\|_{p'}$ and $\|g_{\xi}\|_{\infty}=\sum^{\infty}_{n=1}\|g_{\xi}.\chi_{A_n}\|_{\infty}$. Therefore
\begin{align*}
\sum^{\infty}_{n=1}J_p(A_n)^{\frac{1}{p}}&\leq C\sum^{\infty}_{n=1}\int_{B_{p'}}\|g_{\xi}.\chi_{A_n}\|_{p'}\nu(g_{\xi})\\
&\leq C\int_{B_{p'}}\sum^{\infty}_{n=1}\|g_{\xi}.\chi_{A_n}\|_{p'}\nu(g_{\xi})\\
&=C\int_{B_{p'}}\|g_{\xi}\|_{p'}\nu(g_{\xi})\\
&\leq C\nu(B_{p'})\leq C,
\end{align*}
and similarly $\sum^{\infty}_{n=1}J_1(A_n)\leq C$.
Thus we get the result.
\end{proof}

\begin{thebibliography}{99}
\bibitem[dj]{dj} J. Diestel, H. Jarchow and A. Tonge, Absolutely summing operators. Cambridge Studies in Advanced Mathematics, Camridge University Press, Cambridge, 1995.
\end{thebibliography}
\end{document}
